%======================================================================
% REFERENCIAS TÉCNICAS
%======================================================================

\chapter{Referencias Técnicas}

Este capítulo proporciona información técnica adicional sobre los estándares y tecnologías utilizadas en Sandra UI Consola.

\section{Normativas y Estándares Internacionales}

El desarrollo de Sandra UI Consola está alineado con los marcos más estrictos de la industria:

\subsection{ISO/IEC 27001: Gestión de Seguridad}

La consola implementa un SGSI (Sistema de Gestión de Seguridad de la Información):

\begin{itemize}
\item \textbf{Confidencialidad}: Solo usuarios autorizados acceden a información sensible.
\item \textbf{Integridad}: Protección contra modificaciones no autorizadas.
\item \textbf{Disponibilidad}: El sistema está listo cuando el negocio lo requiere.
\end{itemize}

La certificación ISO 27001 demuestra el compromiso de la organización con la gestión de riesgos de seguridad de la información.

\subsection{ISO 22301: Continuidad del Negocio}

El sistema está diseñado para la Resiliencia Operativa:

\begin{itemize}
\item \textbf{BIA (Business Impact Analysis)}: Identificación de procesos críticos y establecimiento de RTO/RPO.
\item \textbf{Gestión de Respaldos}: Herramientas para generación de dumps y restores.
\item \textbf{Alta Disponibilidad}: Arquitectura distribuida que soporta fallos en nodos.
\end{itemize}

Un RTO (Recovery Time Objective) bajo es crítico para sistemas financieros donde cada minuto de inactividad representa pérdidas significativas.

\subsection{ISO/IEC 25010: Calidad del Software}

\begin{itemize}
\item \textbf{Adecuación Funcional}: Cobertura completa de necesidades operativas.
\item \textbf{Eficiencia de Desempeño}: Tiempos de respuesta ultrarrápidos con Rust y Go.
\item \textbf{Mantenibilidad}: Código limpio y modular para evoluciones constantes.
\end{itemize}

La calidad del software se mide no solo en funcionalidad, sino también en métricas como mantenibilidad, eficiencia y portabilidad.

\section{Endpoints API}

La API de Sandra UI proporciona los siguientes endpoints principales:

\begin{itemize}
    \item \textbf{GET /api/status}: Estado del sistema
    \item \textbf{GET /api/metrics}: Métricas en tiempo real
    \item \textbf{POST /api/auth/login}: Inicio de sesión
    \item \textbf{GET /api/users}: Listar usuarios
    \item \textbf{POST /api/users}: Crear usuario
    \item \textbf{GET /api/logs}: Obtener logs
\end{itemize}

\section{Configuración JSON}

Ejemplo de configuración de conexión:

\begin{verbatim}
{
  "server": {
    "host": "localhost",
    "port": 8080,
    "protocol": "https",
    "timeout": 30000
  },
  "security": {
    "mfa_enabled": true,
    "session_timeout": 3600
  },
  "notifications": {
    "email_enabled": true,
    "slack_webhook": "https://..."
  }
}
\end{verbatim}

\clearpage
